
\documentclass{article}
\usepackage{colortbl}
\usepackage{makecell}
\usepackage{multirow}
\usepackage{supertabular}

\begin{document}

\newcounter{utterance}

\centering \large Interaction Transcript for game `mmlu\_pro', experiment `mmlu\_pro\_engineering', episode 65 with BSU{-}SLIM\_\_Qwen3.5{-}9B{-}Base\_\_playpen{-}prm{-}9b{-}best\_of\_n.
\vspace{24pt}

{ \footnotesize  \setcounter{utterance}{1}
\setlength{\tabcolsep}{0pt}
\begin{supertabular}{c@{$\;$}|p{.15\linewidth}@{}p{.15\linewidth}p{.15\linewidth}p{.15\linewidth}p{.15\linewidth}p{.15\linewidth}}
    \# & \multicolumn{2}{c}{Player} && \multicolumn{2}{c}{Game Master} \\
    \hline

    \theutterance \stepcounter{utterance}  
    & & & \multicolumn{4}{p{0.6\linewidth}}{
        \cellcolor[rgb]{0.9,0.9,0.9}{
            \makecell[{{p{\linewidth}}}]{
                \texttt{\tiny{[P1$\langle$GM]}}
                \texttt{The following are multiple choice questions (with answers) about engineering. Answer only with "The answer is (X)" where X is the correct letter choice.} \\
\\ 
\texttt{Question:} \\
\texttt{A point pole has a strength of 4π * 10\^{}{-}4 weber. The force in newtons on a point pole of 4π * 1.5 * 10\^{}{-}4 weber placed at a distance of 10 cm from it will be} \\
\texttt{Options:} \\
\texttt{A. 25 N.} \\
\texttt{B. 12.5 N.} \\
\texttt{C. 15 N.} \\
\texttt{D. 10 N.} \\
\texttt{E. 7.5 N.} \\
\texttt{F. 20 N.} \\
\texttt{G. 5 N.} \\
\texttt{H. 3.75 N.} \\
\texttt{I. N/A} \\
\texttt{J. N/A} \\
\texttt{The answer is (C)} \\
\\ 
\texttt{Question:} \\
\texttt{The coil of a moving coil meter has 100 turns, is 40 mm long and 30 mm wide. The control torque is 240*10{-}6 N{-}m on full scale. If magnetic flux density is 1Wb/m2 range of meter is} \\
\texttt{Options:} \\
\texttt{A. 2 mA.} \\
\texttt{B. 5 mA.} \\
\texttt{C. 1.5 mA.} \\
\texttt{D. 0.5 mA.} \\
\texttt{E. 6 mA.} \\
\texttt{F. 4 mA.} \\
\texttt{G. 3 mA.} \\
\texttt{H. 1 mA.} \\
\texttt{I. 2.5 mA.} \\
\texttt{J. 3.5 mA.} \\
\texttt{The answer is (A)} \\
\\ 
\texttt{Question:} \\
\texttt{In an SR latch built from NOR gates, which condition is not allowed} \\
\texttt{Options:} \\
\texttt{A. S=0, R=2} \\
\texttt{B. S=2, R=2} \\
\texttt{C. S=1, R=1} \\
\texttt{D. S=1, R={-}1} \\
\texttt{E. S=1, R=2} \\
\texttt{F. S=0, R=0} \\
\texttt{G. S=2, R=0} \\
\texttt{H. S=1, R=0} \\
\texttt{I. S=2, R=1} \\
\texttt{J. S=0, R=1} \\
\texttt{The answer is (C)} \\
\\ 
\texttt{Question:} \\
\texttt{Two long parallel conductors carry 100 A. If the conductors are separated by 20 mm, the force per meter of length of each conductor will be} \\
\texttt{Options:} \\
\texttt{A. 100 N.} \\
\texttt{B. 0.01 N.} \\
\texttt{C. 10 N.} \\
\texttt{D. 0.5 N.} \\
\texttt{E. 0.001 N.} \\
\texttt{F. 5 N.} \\
\texttt{G. 1 N.} \\
\texttt{H. 50 N.} \\
\texttt{I. 0.1 N.} \\
\texttt{J. 0.0001 N.} \\
\texttt{The answer is (I)} \\
\\ 
\texttt{Question:} \\
\texttt{In a 2 pole lap winding dc machine , the resistance of one conductor is 2Ω and total number of conductors is 100. Find the total resistance} \\
\texttt{Options:} \\
\texttt{A. 50Ω} \\
\texttt{B. 1Ω} \\
\texttt{C. 25Ω} \\
\texttt{D. 200Ω} \\
\texttt{E. 10Ω} \\
\texttt{F. 100Ω} \\
\texttt{G. 500Ω} \\
\texttt{H. 150Ω} \\
\texttt{I. 75Ω} \\
\texttt{J. 20Ω} \\
\texttt{The answer is (A)} \\
\\ 
\texttt{Question:} \\
\texttt{Compute the heat gained by the room at 80°F from a vertical plate3 ft. wide and 4 ft. high maintained at 115°F.} \\
\texttt{Options:} \\
\texttt{A. 550.3 Btu/hr} \\
\texttt{B. 684.2 Btu/hr} \\
\texttt{C. 596.4 Btu/hr} \\
\texttt{D. 712.6 Btu/hr} \\
\texttt{E. 739.1 Btu/hr} \\
\texttt{F. 678.9 Btu/hr} \\
\texttt{G. 503.8 Btu/hr} \\
\texttt{H. 625.7 Btu/hr} \\
\texttt{I. 485.2 Btu/hr} \\
\texttt{J. 460.5 Btu/hr} \\
            }
        }
    }
     \\ \\

    \theutterance \stepcounter{utterance}  
    & \multicolumn{4}{p{0.6\linewidth}}{
        \cellcolor[rgb]{0.561,0.737,0.561}{
            \makecell[{{p{\linewidth}}}]{
                \texttt{\tiny{[P1$\rangle$GM]}}
                \texttt{To solve this problem, we calculate the rate of heat transfer by radiation from the vertical plate to the room.} \\
\\ 
\texttt{**1. Identify the parameters:**} \\
\texttt{*   Temperature of the plate (\$T\_p\$): \$115\^{}\textbackslash{}circ\textbackslash{}text\{F\}\$} \\
\texttt{*   Temperature of the room (\$T\_r\$): \$80\^{}\textbackslash{}circ\textbackslash{}text\{F\}\$} \\
\texttt{*   Width of the plate (\$W\$): \$3\textbackslash{}text\{ ft\}\$} \\
\texttt{*   Height of the plate (\$H\$): \$4\textbackslash{}text\{ ft\}\$} \\
\texttt{*   Surface area (\$A\$): \$3\textbackslash{}text\{ ft\} \textbackslash{}times 4\textbackslash{}text\{ ft\} = 12\textbackslash{}text\{ ft\}\^{}2\$} \\
\texttt{*   Stefan{-}Boltzmann constant (\$\textbackslash{}sigma\$): \$5.67 \textbackslash{}times 10\^{}\{{-}8\}\textbackslash{}text\{ W\}/(\textbackslash{}text\{m\}\^{}2\textbackslash{}cdot\textbackslash{}text\{K\}\^{}4)\$ or \$0.1714 \textbackslash{}times 10\^{}\{{-}8\}\textbackslash{}text\{ Btu\}/(\textbackslash{}text\{hr\}\textbackslash{}cdot\textbackslash{}text\{ft\}\^{}2\textbackslash{}cdot\^{}\textbackslash{}circ\textbackslash{}text\{R\}\^{}4)\$. Using English units is more direct here. Let's use \$E = 0.171 \textbackslash{}times 10\^{}\{{-}8\}\textbackslash{}text\{ Btu\}/(\textbackslash{}text\{hr\}\textbackslash{}cdot\textbackslash{}text\{ft\}\^{}2\textbackslash{}cdot\^{}\textbackslash{}circ\textbackslash{}text\{R\}\^{}4)\$ for a black body, but real surfaces have an emissivity (\$\textbackslash{}epsilon\$). Since emissivity is not provided, we assume a standard emissivity for such problems (often around 0.9 or 1.0 for an ideal black body) or that the question implies a specific standard value leads to one of the options. Let's test with \$\textbackslash{}epsilon = 0.9\$ (common for painted metal) and \$\textbackslash{}epsilon = 1.0\$.} \\
\\ 
\texttt{**2. Convert temperatures to Rankine (\$\^{}\textbackslash{}circ\textbackslash{}text\{R\}\$):**} \\
\texttt{*   \$T\_p (\^{}\textbackslash{}circ\textbackslash{}text\{R\}) = 115 + 459.67 = 574.67\^{}\textbackslash{}circ\textbackslash{}text\{R\}\$} \\
\texttt{*   \$T\_r (\^{}\textbackslash{}circ\textbackslash{}text\{R\}) = 80 + 459.67 = 539.67\^{}\textbackslash{}circ\textbackslash{}text\{R\}\$} \\
\texttt{*   Alternatively, using the approximation \$\^{}\textbackslash{}circ\textbackslash{}text\{R\} = \^{}\textbackslash{}circ\textbackslash{}text\{F\} + 460\$:} \\
\texttt{    *   \$T\_p \textbackslash{}approx 575\^{}\textbackslash{}circ\textbackslash{}text\{R\}\$} \\
\texttt{    *   \$T\_r \textbackslash{}approx 540\^{}\textbackslash{}circ\textbackslash{}text\{R\}\$} \\
\\ 
\texttt{**3. Calculate temperature difference:**} \\
\texttt{*   \$T\_p\^{}4 {-} T\_r\^{}4 = 575\^{}4 {-} 540\^{}4\$} \\
\texttt{*   \$575\^{}4 = 1,107,464,062,500\$} \\
\texttt{*   \$540\^{}4 = 850,305,600,000\$} \\
\texttt{*   Difference = \$257,158,462,500\$} \\
\\ 
\texttt{**4. Apply the heat transfer formula:**} \\
\texttt{\$Q = \textbackslash{}epsilon \textbackslash{}cdot \textbackslash{}sigma \textbackslash{}cdot A \textbackslash{}cdot (T\_p\^{}4 {-} T\_r\^{}4)\$} \\
\\ 
\texttt{Using the standard English unit constant \$\textbackslash{}sigma \textbackslash{}approx 1.714 \textbackslash{}times 10\^{}\{{-}9\}\textbackslash{}text\{ Btu\}/(\textbackslash{}text\{hr\}\textbackslash{}cdot\textbackslash{}text\{ft\}\^{}2\textbackslash{}cdot\^{}\textbackslash{}circ\textbackslash{}text\{R\}\^{}4)\$ (Note: \$0.1714 \textbackslash{}times 10\^{}\{{-}8\}\$ is the same thing).} \\
\\ 
\texttt{Let's try calculating without emissivity (assuming black body, \$\textbackslash{}epsilon=1\$) first to see if it matches an option.} \\
\texttt{\$Q = 1 \textbackslash{}cdot (1.714 \textbackslash{}times 10\^{}\{{-}9\}) \textbackslash{}cdot 12 \textbackslash{}cdot (2.5716 \textbackslash{}times 10\^{}\{11\})\$} \\
\texttt{\$Q = 1.714 \textbackslash{}times 12 \textbackslash{}times 257.16\$} \\
\texttt{\$Q \textbackslash{}approx 20.568 \textbackslash{}times 257.16 \textbackslash{}approx 5284\textbackslash{}text\{ Btu/hr\}\$} \\
\\ 
\texttt{This result is nowhere near the options (which are in the range of 460–740). This suggests the units might be in **Watts** and the options are mislabeled, or the temperatures are in **Kelvin**, or the formula involves convection, or the surface area is different.} \\
\\ 
\texttt{**Alternative Interpretation: Convection**} \\
\texttt{If the heat transfer is by convection (Newton's Law of Cooling):} \\
\texttt{\$Q =} \\
            }
        }
    }
    & & \\ \\

    \theutterance \stepcounter{utterance}  
    & & & \multicolumn{2}{p{0.3\linewidth}}{
        \cellcolor[rgb]{0.9,0.9,0.9}{
            \makecell[{{p{\linewidth}}}]{
                \texttt{\tiny{[GM$|$GM]}}
                \texttt{ParseError: Response doesn't contain any answer, but the expected format give one.} \\
            }
        }
    }
    & & \\ \\

    \theutterance \stepcounter{utterance}  
    & & & \multicolumn{2}{p{0.3\linewidth}}{
        \cellcolor[rgb]{0.9,0.9,0.9}{
            \makecell[{{p{\linewidth}}}]{
                \texttt{\tiny{[GM$|$GM]}}
                \texttt{game\_result = ABORT} \\
            }
        }
    }
    & & \\ \\

\end{supertabular}
}

\end{document}
